\chapter{Introduction} \label{cap:introducao}

%Este documento é baseado no exemplo de referência de uso da classe
%\textsf{abntex2} e do pacote \textsf{abntex2cite}.
%Este documento é uma referência e os manuais do \abnTeX\ \cite{abntex2classe,abntex2cite,abntex2cite-alf} e da classe \textsf{memoir} \cite{memoir} possuem informações completas.

%(Observe no parágrafo anterior como são feitas as citações! As referências estão no arquivo abntex2-modelo-references.bib, seguindo um formato específico para as informações necessárias. Ah, se você quiser fazer uma citção com o nome do autor no meio do texto, por exemplo dizendo que o \citeonline{abntex2classe} fez um trabalho importante, você usa o citeonline).

%Os capítulo, seções e conteúdos aqui propostos são baseados no documento \textit{Sumário para Monografia de Projeto de Formatura}.

%% CONTEXT
\section{Context}

%%%%%%%%%%%%%%%%%%%%%%%%%%
\section{Motivation}

%\begin{itemize}
%    \item Apresentar o estado de arte resumido do assunto que será o tema do desenvolvimento do trabalho, elaborado com base nos trabalhos consultados.
%    \item Citar os trabalhos consultados no texto.
%    \item Contexto em que o trabalho será desenvolvido.
%\end{itemize}

%%%%%%%%%%%%%%%%%%%%%%%%%%
\section{Objectives}

%\begin{itemize}
%    \item Apresentar o objetivo do trabalho de forma precisa e concisa.
%    \item O que é o trabalho?
%\end{itemize}

%%%%%%%%%%%%%%%%%%%%%%%%%%

%%%%%%%%%%%%%%%%%%%%%%%%%%
\section{Related works}

%\begin{itemize}
%    \item Apresentar porque o trabalho desenvolvido é importante (importância e necessidade para a sociedade, comparação com trabalhos relevantes consultados etc.).
%    \item Citar os trabalhos consultados no texto.
%    \item Por que o trabalho é importante?
%\end{itemize}

%%%%%%%%%%%%%%%%%%%%%%%%%%
\section{Organisation}

%Descrever sucintamente os capítulos e as demais partes do trabalho.
%Lembre-se de colocar rótulos nos capítulos e seções (com label{} ) e depois utilizar o ref{}
%Por exemplo, o Capítulo \ref{cap:conceitos} apresenta conceitos....


